%%%%%%%%%%%%%%%%%%%%%%%%%%%%%%%%%%%%%%%%%%%%%%%%%%%%%%%%%%%%%%%
%% OXFORD THESIS TEMPLATE

% Use this template to produce a standard thesis that meets the Oxford University requirements for DPhil submission
%
% Originally by Keith A. Gillow (gillow@maths.ox.ac.uk), 1997
% Modified by Sam Evans (sam@samuelevansresearch.org), 2007
% Modified by John McManigle (john@oxfordechoes.com), 2015
%
% This version Copyright (c) 2015-2023 John McManigle
%
% Broad permissions are granted to use, modify, and distribute this software
% as specified in the MIT License included in this distribution's LICENSE file.
%

% I've (John) tried to comment this file extensively, so read through it to see how to use the various options.  Remember
% that in LaTeX, any line starting with a % is NOT executed.  Several places below, you have a choice of which line to use
% out of multiple options (eg draft vs final, for PDF vs for binding, etc.)  When you pick one, add a % to the beginning of
% the lines you don't want.


%%%%% CHOOSE PAGE LAYOUT
% The most common choices should be below.  You can also do other things, like replacing "a4paper" with "letterpaper", etc.

% This one will format for two-sided binding (ie left and right pages have mirror margins; blank pages inserted where needed):
\documentclass[a4paper,twoside]{ociamthesis}
% This one will format for one-sided binding (ie left margin > right margin; no extra blank pages):
%\documentclass[a4paper]{ociamthesis}
% This one will format for PDF output (ie equal margins, no extra blank pages):
%\documentclass[a4paper,nobind]{ociamthesis} 

%%%%% SELECT YOUR DRAFT OPTIONS
% Three options going on here; use in any combination.  But remember to turn the first two off before
% generating a PDF to send to the printer!

% This adds a "DRAFT" footer to every normal page.  (The first page of each chapter is not a "normal" page.)
\fancyfoot[C]{\emph{DRAFT Printed on \today}}  

% This highlights (in blue) corrections marked with (for words) \mccorrect{blah} or (for whole
% paragraphs) \begin{mccorrection} . . . \end{mccorrection}.  This can be useful for sending a PDF of
% your corrected thesis to your examiners for review.  Turn it off, and the blue disappears.
\correctionstrue


%%%%% BIBLIOGRAPHY SETUP
% Note that your bibliography will require some tweaking depending on your department, preferred format, etc.
% The options included below are just very basic "sciencey" and "humanitiesey" options to get started.
% If you've not used LaTeX before, I recommend reading a little about biblatex/biber and getting started with it.
% If you're already a LaTeX pro and are used to natbib or something, modify as necessary.
% Either way, you'll have to choose and configure an appropriate bibliography format...

% The science-type option: numerical in-text citation with references in order of appearance.
\usepackage[style=numeric-comp, sorting=none, backend=biber, doi=false, isbn=false]{biblatex}


% \usepackage[style=authoryear, sorting=nyt, backend=biber, maxcitenames=2, useprefix, doi=false, isbn=false]{biblatex}
\newcommand*{\bibtitle}{References}

% This makes the bibliography left-aligned (not 'justified') and slightly smaller font.
\renewcommand*{\bibfont}{\raggedright\small}

% Change this to the name of your .bib file (usually exported from a citation manager like Zotero or EndNote).
\addbibresource{references.bib}


% Uncomment this if you want equation numbers per section (2.3.12), instead of per chapter (2.18):
%\numberwithin{equation}{subsection}



%%%%% THESIS / TITLE PAGE INFORMATION
% Everybody needs to complete the following:
\title{Title}
\author{Neil Natarajan}
\college{New College}

% Master's candidates who require the alternate title page (with candidate number and word count)
% must also un-comment and complete the following three lines:
%\masterssubmissiontrue
%\candidateno{933516}
%\wordcount{28,815}

% Uncomment the following line if your degree also includes exams (eg most masters):
%\renewcommand{\submittedtext}{Submitted in partial completion of the}
% Your full degree name.  (But remember that DPhils aren't "in" anything.  They're just DPhils.)
\degree{Doctor of Philosophy}
% Term and year of submission, or date if your board requires (eg most masters)
\degreedate{Trinity 2024}


%%%%% YOUR OWN PERSONAL MACROS
% This is a good place to dump your own LaTeX macros as they come up.

% To make text superscripts shortcuts
	\renewcommand{\th}{\textsuperscript{th}} % ex: I won 4\th place
	\newcommand{\nd}{\textsuperscript{nd}}
	\renewcommand{\st}{\textsuperscript{st}}
	\newcommand{\rd}{\textsuperscript{rd}}

%%%%% THE ACTUAL DOCUMENT STARTS HERE
\begin{document}



%%%%% CHOOSE YOUR LINE SPACING HERE
% This is the official option.  Use it for your submission copy and library copy:
\setlength{\textbaselineskip}{22pt plus2pt}
% This is closer spacing (about 1.5-spaced) that you might prefer for your personal copies:
%\setlength{\textbaselineskip}{18pt plus2pt minus1pt}

% You can set the spacing here for the roman-numbered pages (acknowledgements, table of contents, etc.)
\setlength{\frontmatterbaselineskip}{17pt plus1pt minus1pt}

% Leave this line alone; it gets things started for the real document.
\setlength{\baselineskip}{\textbaselineskip}


%%%%% CHOOSE YOUR SECTION NUMBERING DEPTH HERE
% You have two choices.  First, how far down are sections numbered?  (Below that, they're named but
% don't get numbers.)  Second, what level of section appears in the table of contents?  These don't have
% to match: you can have numbered sections that don't show up in the ToC, or unnumbered sections that
% do.  Throughout, 0 = chapter; 1 = section; 2 = subsection; 3 = subsubsection, 4 = paragraph...

% The level that gets a number:
\setcounter{secnumdepth}{2}
% The level that shows up in the ToC:
\setcounter{tocdepth}{2}


%%%%% ABSTRACT SEPARATE
% This is used to create the separate, one-page abstract that you are required to hand into the Exam
% Schools.  You can comment it out to generate a PDF for printing or whatnot.
\begin{abstractseparate}
	In a globalised world, scholarship programs face pressures to select the ``best'' applicants fairly across regions, but current low-tech selection processes struggle to meet this burden. To address this, we studied two global scholarship programs, exploring data-driven Decision Support Tools (DSTs) for selection. One program was interested in utilising existing AI tools wherever possible. We explored using existing post-hoc explainable AI (xAI) tools as DSTs in a variety of contexts and using existing detection software as DSTs to understand generative AI (genAI)'s role in essays. We developed the \emph{Decision Matrix} framework, categorising decisions by \emph{stage} (\emph{in-process} or \emph{ex-post}) and \emph{stakes} (\emph{high} or \emph{low}), and applied it to assess both technologies as DSTs. Both proved ineffective for \emph{in-process} decisions and useful \emph{ex-post}.

To design a DST for \emph{in-process} decision-making, we engaged in participatory design with both programs, focusing on diversity considerations in selection. Six prototypes were created, with practitioners showing eagerness to use them across the \emph{Decision Matrix}. To validate this enthusiasm, we implemented one design as a technology probe and evaluated its impact. The selected cohort's diversity and performance improved, demonstrating the tool's ability to support high-stakes \emph{in-process} decisions.

Our findings highlight the need for data-driven and AI-based DSTs across the \emph{Decision Matrix}. We propose \emph{Selection-Oriented AI}, focused on the pro-social goals of selection. We conclude with a call for AI-driven DSTs that balance practitioners' needs while optimising selection outcomes for social benefit. % Create an abstract.tex file in the 'text' folder for your abstract.
\end{abstractseparate}


% JEM: Pages are roman numbered from here, though page numbers are invisible until ToC.  This is in
% keeping with most typesetting conventions.
\begin{romanpages}

% JEM: By default, this template uses the traditional Oxford "Belt Crest". Un-comment the following
% line to use the newer, "Blue Square" logo:
% \renewcommand{\crest}{{\includegraphics[width=4.2cm, height=4.2cm]{figures/newlogo.pdf}}}

% Title page is created here
\maketitle

%%%%% DEDICATION -- If you'd like one, un-comment the following.
%\begin{dedication}
%This thesis is dedicated to\\
%someone\\
%for some special reason\\
%\end{dedication}

%%%%% ACKNOWLEDGEMENTS -- Nothing to do here except comment out if you don't want it.
\begin{acknowledgements}
 	I would like to thank my advisors Reuben Binns and Nigel Shadbolt: Reuben, for his constant dedication to his students and willingness to run with any idea I came up with; Nigel, for his clear framing and pointed critiques. I would also like to thank my colleagues in Oxford Computer Science's HCAI theme, with special thanks to Sruthi, Ulrik, Jake, and Thomas. I would like to thank my family for their support and my friends for my sanity.
\end{acknowledgements}

%%%%% ABSTRACT -- Nothing to do here except comment out if you don't want it.
\begin{abstract}
	In a globalised world, scholarship programs face pressures to select the ``best'' applicants fairly across regions, but current low-tech selection processes struggle to meet this burden. To address this, we studied two global scholarship programs, exploring data-driven Decision Support Tools (DSTs) for selection. One program was interested in utilising existing AI tools wherever possible. We explored using existing post-hoc explainable AI (xAI) tools as DSTs in a variety of contexts and using existing detection software as DSTs to understand generative AI (genAI)'s role in essays. We developed the \emph{Decision Matrix} framework, categorising decisions by \emph{stage} (\emph{in-process} or \emph{ex-post}) and \emph{stakes} (\emph{high} or \emph{low}), and applied it to assess both technologies as DSTs. Both proved ineffective for \emph{in-process} decisions and useful \emph{ex-post}.

To design a DST for \emph{in-process} decision-making, we engaged in participatory design with both programs, focusing on diversity considerations in selection. Six prototypes were created, with practitioners showing eagerness to use them across the \emph{Decision Matrix}. To validate this enthusiasm, we implemented one design as a technology probe and evaluated its impact. The selected cohort's diversity and performance improved, demonstrating the tool's ability to support high-stakes \emph{in-process} decisions.

Our findings highlight the need for data-driven and AI-based DSTs across the \emph{Decision Matrix}. We propose \emph{Selection-Oriented AI}, focused on the pro-social goals of selection. We conclude with a call for AI-driven DSTs that balance practitioners' needs while optimising selection outcomes for social benefit.
\end{abstract}

%%%%% MINI TABLES
% This lays the groundwork for per-chapter, mini tables of contents.  Comment the following line
% (and remove \minitoc from the chapter files) if you don't want this.  Un-comment either of the
% next two lines if you want a per-chapter list of figures or tables.
\dominitoc % include a mini table of contents
%\dominilof  % include a mini list of figures
%\dominilot  % include a mini list of tables

% This aligns the bottom of the text of each page.  It generally makes things look better.
\flushbottom

% This is where the whole-document ToC appears:
\tableofcontents

% Uncomment to generate a list of figures:
% \listoffigures
% 	\mtcaddchapter

% \mtcaddchapter is needed when adding a non-chapter (but chapter-like) entity to avoid confusing minitoc

% Uncomment to generate a list of tables:
%\listoftables
%	\mtcaddchapter

%%%%% LIST OF ABBREVIATIONS
% This example includes a list of abbreviations.  Look at text/abbreviations.tex to see how that file is
% formatted.  The template can handle any kind of list though, so this might be a good place for a
% glossary, etc.
% First parameter can be changed eg to "Glossary" or something.
% Second parameter is the max length of bold terms.
\begin{mclistof}{List of Abbreviations}{3.2cm}
    \item[AI] Artificial Intelligence
    \item[xAI] Explainable Artificial Intelligence
    \item[DEI] Diversity, Equity, and Inclusion.
    \item[IAI] Interpretable Artificial Intelligence
    \item[DST] Decision Support Tool
    \item[HitL] Human in the Loop 
    \item[IAIDST] Interpretable, Artificially Intelligent Decision Support Tool
    \item[TI] Talent Identification 
\end{mclistof}
See Appendix \ref{app:terminology} for definitions of terms used in this Thesis, including these abbreviations.

% The Roman pages, like the Roman Empire, must come to its inevitable close.
\end{romanpages}


%%%%% CHAPTERS
% Add or remove any chapters you'd like here, by file name (excluding '.tex'):
\flushbottom
% \begin{savequote}[8cm]
% \textlatin{Neque porro quisquam est qui dolorem ipsum quia dolor sit amet, consectetur, adipisci velit...}

% There is no one who loves pain itself, who seeks after it and wants to have it, simply because it is pain...
%   \qauthor{--- Cicero's \textit{de Finibus Bonorum et Malorum}}
% \end{savequote}

\chapter{\label{ch:intro}Introduction} 

\minitoc

Since early theorising about Artificial Intelligence (AI), experts have noted the transformative potential of such technology. Now, with the popular advent of Machine Learning (ML) models and transformer architecture in particular, we see rapid mass-market adoption of a variety of AI tools. In many fields which deal with sensitive decision-making, this adoption introduces new ethical challenges. In the fields of talent and Talent Identification (TI) in particular, many of these ethical challenges centre on decision tool interpretability; diversity, equity, and inclusion (DEI); and applicant usage of generative AI. In this section, we lay out the case for the application of AI systems (specifically interpretable AI decision assistants) to these issues in TI, paying special attention to the potential pitfalls and broader implications of our proposal.

\section{A Brief Account of Talent Identification, its Adoption of Artificial Intelligence, and Challenges Arising}
(To-do)
% Intro needs a sociological account of talent identification. What is it? How does it differ from recruitment? What niche within TI do I focus on? I talk about this a bit in the organisation discussion but should go into greater detail here.

\section{Motivation}
(To-do) This is where I explain why I'm doing this research. I explain why care about talent identification and how IAIDSTs can solve problems here. I respond to the notion that we should simply not use AI here, and highlight problems in the field already. I explain why this is a problem for Computer Science as well as for Talent Science (tl;dr: additional reasons to care about things xai overlooks because of the application domain)

\section{Partner Organisations}
(To-do) Describe the two organisations partnered with for this research.

\section{Contributions}
In this thesis, we make a number of theoretical, technical and practical contributions. These are:

\begin{itemize}
    \item We articulate where talent identification experts should use IAIDSTs to assist in their decision-making, and which IAI systems are appropriate for which purposes
    \item We apply research-through-design methods to talent identification to produce a novel approach to the creation of human-centric selection tooling
    \item We develop a novel method for visualising and understanding cohort-level decisions surrounding diversity called the Selection Possibility Frontier
    \item We further implement a novel algorithm for approximating the Selection Possibility Frontier
    \item We produce several qualitative, quantitative, and mixed-modal datasets relating to talent identification
    \item We demonstrate a yet-undemonstrated risk associated with naive application of post-hoc explainability to a range of problems
    \item We demonstrate the utility and limitation of detection tools for generative-AI-written essays in a scholarship context, and explore implications for scholarship programs
\end{itemize}

\section{Thesis Structure}
In this chapter, we explore our motivation for focusing on the application of IAIDSTs to talent identification, enumerate our contributions, and provide an account of challenges talent identification. Among these challenges, we focus the remainder of this thesis on three: limited human understanding of the AI tools, limited ability to consider diversity, and the risks of applicant usage of generative AI.

In Chapter \ref{ch:background}, we provide a background on the field of talent identification, the field of explainable AI, and the field of diversity, equity, and inclusion; we further explore intersections between these fields of research and identify gaps in the literature. Chapters \ref{ch:xailimitations} and \ref{ch:usingxai} explore our first challenge: limited human understanding of the AI tools. In Chapter \ref{ch:xailimitations}, we explore the limitations of popular xAI methods in the context of talent identification. In Chapter \ref{ch:usingxai}, we explore the use of SHAP explanations in the context of talent identification. Chapters \ref{ch:iaicasestudy} and \ref{ch:spf}  address challenges related to consideration of diversity in TI. In Chapter \ref{ch:iaicasestudy}, we present a participatory design study, which aims to use IAIDSTs to help TI professionals account for diversity. In Chapter \ref{ch:spf} we implement one of the designs from Chapter \ref{ch:iaicasestudy}. Chapter \ref{ch:studentuse} attempts to mitigate the risks of applicant usage of generative AI. It explores the use of generative AI detection tools as DSTs in the context of talent identification. Finally, in Chapter \ref{ch:discussion}, we discuss our design recommendations, note limitations of our approach, and conclude.

\section{Publications}

\begin{itemize}
    \item \textcite{natarajan_trust_2023} Trust Explanations to do What They Say
    \item \textcite{...} `You Trust Me, But Should You?': Misleading Explanations of AI Outputs
\end{itemize}
% \begin{savequote}[8cm]
% \textlatin{Neque porro quisquam est qui dolorem ipsum quia dolor sit amet, consectetur, adipisci velit...}

% There is no one who loves pain itself, who seeks after it and wants to have it, simply because it is pain...
%   \qauthor{--- Cicero's \textit{de Finibus Bonorum et Malorum}}
% \end{savequote}

\chapter{\label{ch:2-methods}Methods} 

\minitoc

\section{To-do}
To-do % Where is background? Is it part of methods? Do I call it "prior literature and methods?"
% \begin{savequote}[8cm]
% Alles Gescheite ist schon gedacht worden.\\
% Man muss nur versuchen, es noch einmal zu denken.

% All intelligent thoughts have already been thought;\\
% what is necessary is only to try to think them again.
%   \qauthor{--- Johann Wolfgang von Goethe \cite{von_goethe_wilhelm_1829}}
% \end{savequote}

\chapter{\label{ch:3-xailimitations}Limitations of Post-Hoc Explainable AI}

\minitoc

\section{Introduction}
Increasingly, decisions affecting the lives of lay people are made by AI algorithms. And while these algorithms may be useful, they can also be dangerous. Users of AI systems desire an understanding of how these systems functions and why they yield the outputs they do, so that they may respond appropriately to the outputs \cite{binns_human_2022}. When AI systems make decisions impacting the user, user insight into the decision-making process allows for recourse \cite{ustun_actionable_2019}. However, when decision-making systems involve both an AI and human component, user insight into AI outputs is crucial in both (1) inducing appropriate reliance or scepticism in the AI as warranted and (2) optimising the overall performance of the decision-making system. This is especially true in high-stakes domains such as healthcare, finance, and criminal justice, where the decisions made by AI systems can have significant consequences for individuals \cite{binns_human_2022,ustun_actionable_2019,wachter_counterfactual_2017}.

The growing field of Explainable Artificial Intelligence (xAI) aims to develop methods for explaining algorithms. While some models are intrinsically interpretable \cite{rudin_interpretable_2021}, many more complicated machine learning algorithms are difficult to interpret. In order to explain these more complex models, various dedicated explanation algorithms have been proposed. Within the field of xAI, post-hoc, model-agnostic explainability methods remain popular due to their flexibility and ease-of-application \cite{molnar_interpretable_2019}. Indeed, the ability of methods like SHAP, LIME, and Scoped Anchors to treat any model as a black-box allows for applications of explainability to models that would be otherwise inscrutable \cite{lundberg_unified_2017,ribeiro_why_2016,ribeiro_anchors_2018}. However, there are several known limitations that can impact their utility as IAIDSTs; in some cases, these limitations undermine the purpose of explanation, or even to deceive and mislead users \cite{wang_transparency_2022}. 

In this Chapter, we explore a disconnect between user trust in an AI system and that system's trustworthiness as a result of post-hoc explainable AI tools. We begin by citing a series of papers yielding this result in a variety of contexts, and proceed to present our own work confirming that this disconnect persists in talent identification tasks. We then identify theories of when and why this disconnect occurs, and propose a series of experiments to test these theories. 

In Chapter \ref{ch:4-xaicasestudy}, we implement these experiments and present the results.

\section{Trust and Trustworthiness in Explainable AI}
User trust is central to the utility and limitations of any user-facing system. In the context of IAIDSTs, user trust in the system facilitates the system's adoption and use. Explainability is crucial here, as an oft-cited goal of transparency is increasing user trust in AI systems \cite{wang_transparency_2022}. However, though trust is much-discussed, trustworthiness is often overlooked.

\textcite{jacovi_formalizing_2021} set out rules for trust in an algorithmic context. They set out a basic philosophical analysis of trust (in the general, non-algorithmic sense): A trusts B if and only if A believes that B will act in A's best interest and accepts vulnerability to B's actions. Trust often has limited scope; typically, A will trust B regarding some particular actions or motivations, but not others.

In the context of trust in AI, this definition is only slightly changed. \textcite{jacovi_formalizing_2021} characterise trust in an AI system by two properties: ``the vulnerability of the user, and the ability to anticipate the impact of the AI model's decisions''; \textcite{vereschak_how_2021} similarly isolate three elements: ``trust is linked to a situation of vulnerability and positive expectations, and is an attitude''; \textcite{lee_trust_2004} give a similar definition of trust: ``An attitude that an agent will achieve an individual's goal in a situation characterized by uncertainty and vulnerability''. In all definitions, we see \emph{vulnerability} emerge as a key concept, and we variably also see that trust is characterised by \emph{uncertainty} and \emph{expectations}.

An important distinction here can be drawn between whether someone or something is trusted and whether that trust is well-placed; i.e. it is worthy of trust \cite{hardin_trust_2002}. In the machine context, Jacovi et al. argue, an algorithm is worthy of trust if and only if there exists some contract that the algorithm promises (or that the algorithm's creators and implementors promise) to uphold. They term this a `contractual' model of trust. We adopt the nomenclature of contractual trust and use it to frame breakdowns in trust in specific post-hoc explanations.

\section{A Series of Concerning Results}
\subsection{The Dangers of Transparency}
It is often assumed in the xAI literature that transparency of AI systems is always desirable \cite{molnar_interpretable_2019,miller_explanation_2017-1}. However, \textcite{wang_transparency_2022} dispel this claim. In particular, they demonstrate how, rather than empowering users and stakeholders, some implementations of algorithmic transparency serve the system's deployers, rather than its users, either through a false sense of understanding or the enforcement of norms and power structures. They examine the USA's FICO credit score, which releases the factors that go into the score, the data sources for these factors, and general guidelines on how these factors are used. They find specific information on the use of these factors lacking, and conclude that, rather than empowering users to debate the ethics of certain decisions, this transparency only serves to enforce certain behaviours among the users \cite{wang_transparency_2022}.

\textcite{ustun_actionable_2019} note a similar danger in transparency. They argue that explanations of AI systems making decisions about at end-users should empower those end users to alter the system's determinations. In particular, they introduce the concept of `Actionable Recourse', which consists of a specific set of actions an end-user can take to change the AI's determination, and demonstrate how to calculate these for linear models. Key to actionable recourse is that the actions an explainee should take in response to the explanation are clear in the explanation \cite{ustun_actionable_2019}.

We extend this concept to the context of IAIDSTs. Explanations aimed at decision-makers need not provide them \emph{recourse}, as decision-makers already posses the ability to override the AI system. However, they should still provide \emph{actionable} information, so that the explanations should contain information relevant to the decision-makers choice of what to do with the machine recommendation.

\subsection{Black Box Explanations Increase User Trust in AI Systems}
A number of studies exploring popular post-hoc, black-box explanation algorithms have found that these methods tend to increase overall user trust in the system being explained \cite{ford_play_2020,jacobs_how_2021,wang_transparency_2022}. 

\textcite{ford_play_2020} explore this effect in the context of a machine decision-maker with a human evaluator. They run a study examining the impact of post-hoc explanations-by-example and error-rates on people's perceptions of a black-box classifier classifying images from the MNIST dataset. They show that presenting `case-based explanations' (a series of three important data points in the training of the model; elsewhere, we term this an `influential instances' explanation and use `case-based' more broadly) lead participants to perceive miss-classifications as more correct \cite{ford_play_2020}.

\textcite{jacobs_how_2021} extend this result to a human-in-the-loop context. They study the effect of machine recommendation on a clinician's ability to select antidepressant treatments, and find that providing clinicians an incorrect algorithmically generated treatment list lowers the accuracy of clinicians' own treatment lists. Notably, they do not isolate this result to the presence of an explanation, but rather demonstrate how two explainable AI systems harm clinician antidepressant selection relative to a placebo group. However, they also find that feature-based explanations mislead clinicians more than heuristic-based explanations. Modern best practices outlined by \textcite{miller_explanation_2017-1} prefer the usage of feature-based explanations, so we find this particularly alarming \cite{jacobs_how_2021}.

\textcite{mccradden_when_2021} responds to \textcite{jacobs_how_2021} questioning the role of accuracy in explanations of IAIDSTs. They note the focus of \textcite{jacobs_how_2021} on improving \emph{accuracy} of the treatment lists, and aim to do this with IAIDSTs. However, though these systems may optimise for accuracy, they notably fail (at least in these specific instances) to facilitate clinicians' ability to \emph{help patients}, which is a clinicians' primary goal \cite{mccradden_when_2021}.

In a broader context, while `helping patients' doesn't describe the goal of all IAIDSTs, neither does `improving accuracy'. And, though increasing decision-maker trust regardless of system veracity may improve overall accuracy, it may hamper the broader goal of the decision-making system.

\subsection{Impacts on Talent Identification}
It should be noted that these effects are not universal. \textcite{mohseni_trust_nodate} measure user trust in an AI fake news detection system over time, and find the profile of the users to be just as important as the type of explanation. Though they demonstrate that different explanatory conditions have differential effects on which profile participants are likely to exhibit, they also cluster trust into five profiles by user: consistent over-trust, consistent under-trust, consistent trust, trust gains continually, and trust decreases continually. These profiles suggest that the same explanations will impact different groups performing different tasks differently. Thus, while an IAIDST may lead clinicians selecting antidepressants astray, it may have no effect (or even the opposite effect) on talent identification tasks \cite{mohseni_trust_nodate}.

In the next section, we undertake a series of experiments to test the implications of these limitations for talent identification.

\section{Misleading Explanations of AI Outputs in Talent Identification}
To-do: \emph{Explainable AI (xAI) methods are often motivated by the need to increase user trust in AI systems. However, increased trust may not always be desirable. While too little trust might lead people to neglect AI systems when they are correct, too much trust can also be a risk. Explanations which encourage misplaced trust in incorrect AI outputs might mislead humans into making bad decisions. We investigate the extent to which three different xAI methods create too much trust in their underlying AI systems. Examining Shapley-based Additive Explanations (SHAP), Scoped Anchors, and a confidence-based explanation that presents the model's confidence in an output, we conduct experiments where participants perform tasks with xAI assistance. We find that when the system is wrong, SHAP or Confidence explanations still increase trust. Scoped Anchors explanations, by contrast, increase participants' confidence in their own estimations regardless of the system's correctness. We discuss implications for the design and deployment of xAI in HitL tasks.}

\section{Trust Explanations to Do What They Say} % Is this part of intro or methods? It's almost more of a commentary on the state of the literature, so possibly fits better with background.
To-do: \emph{Trusting an algorithm without cause may lead to abuse, and mistrusting it may similarly lead to disuse. Trust in an AI is only desirable if it is warranted; thus, calibrating trust is critical to ensuring appropriate use. In the name of calibrating trust appropriately, AI developers should provide contracts specifying use cases in which an algorithm can and cannot be trusted. Automated explanation of AI outputs is often touted as a method by which trust can be built in the algorithm. However, automated explanations arise from algorithms themselves, so trust in these explanations is similarly only desirable if it is warranted. Developers of algorithms explaining AI outputs (xAI algorithms) should provide similar contracts, which should specify use cases in which an explanation can and cannot be trusted.}

% Clearly, post-hoc and intrinsically interpretable are very different. We clearly need to study them separately. That's what the next two chapters do
% \begin{savequote}[8cm]
% Alles Gescheite ist schon gedacht worden.\\
% Man muss nur versuchen, es noch einmal zu denken.

% All intelligent thoughts have already been thought;\\
% what is necessary is only to try to think them again.
%   \qauthor{--- Johann Wolfgang von Goethe \cite{von_goethe_wilhelm_1829}}
% \end{savequote}

\chapter{\label{ch:4-xaicasestudy}Using Post-Hoc Explainable AI to Identify Talent}

\minitoc

\section{Introduction}
To-do

\section{Explaining Applicant Resumes with Highlighting and Blinding}
We develop a white-box explanation algorithm for highlighting the resumes of job applicants and scoring them. We use only the highlights themselves, opting for an evaluative AI approach, and provide a contract for the intended use case of our algorithm. We then implement this tool in the hiring process of an organisation and conduct a field study on its impact on hiring decision-making. We discover...

(To-do)

\section{Explaining an Applicant Scoring Algorithm with Shapley Values}
We develop a SHAP-based procedure for explaining the decisions made by an applicant scoring algorithm used in a talent investment organisation. We take into account the original functions of SHAP explanations and provide a contract for the intended use case of our procedure: to systematically reveal insights about the underlying algorithm itself, help us better understand counterintuitive results, and instruct us on where this algorithm should be modified. We explore decisions made by the talent investment organisation and analyse the applicant scoring algorithm using the SHAP tool. We discover... Finally, we discuss the potential for extending this tool to use in decision-making and note strategies to combat mis-calibrated trust.

(To-do)

\section{Detecting Generative AI Usage in Application Essays} % Slightly worried about this. It's narrative fit with ch4 isn't as strong as the others. It's not a cohort-type issue here. It's also not "identifying talent" in the way the other two are. Does this get its own chapter? Maybe this goes after the IAI bit. Then it's something like "we've talked about TI people using AI, but an obvious shift in TI is being caused by students using AI. We need to understand how to detect this / what it means. The student use here might threaten everything else we've done, and we need to scope the problem."
%Alt: if xAI filtering and CV screening becomes the norm, can applicants game this? Can you produce outputs such that the determinations favour you in some way?
Student use of generative AI in essay-writing creates new challenges for education in the marking of essays and essay-based selection for scholarships, fellowships, and universities. In theory, software purporting to detect AI-generated content can act as a post-hoc explanation of generated content, and thus offers a plausible solution to educators inquiring about the use of generative AI in student essays. In practice, little research has attempted to validate such solutions in real-world situations or to examine their limitations. We present an empirical case study exploring the efficacy and implications of using one such detection product, GPTZero, in the selection process for a program looking for talented young people from around the world. We observe that GPTZero does not perform sufficiently well for the program to disqualify applicants on its basis alone. Also, GPTZero's scores are heterogeneously biased across geographical and gender groups. However, we find GPTZero accurate enough to conduct useful aggregate analyses of potential vulnerability to AI-enabled attacks on the integrity of the program's application process. We find evidence for only limited use of AI-generated text in the program's latest application cycle and no evidence partner organizations used AI-generated text at scale to defraud the program's referral program.

(To-do)

% \begin{savequote}[8cm]
% Alles Gescheite ist schon gedacht worden.\\
% Man muss nur versuchen, es noch einmal zu denken.

% All intelligent thoughts have already been thought;\\
% what is necessary is only to try to think them again.
%   \qauthor{--- Johann Wolfgang von Goethe \cite{von_goethe_wilhelm_1829}}
% \end{savequote}

\chapter{\label{ch:5-iaicasestudy}Using Intrinsically Interpretable Models to Identify Talent} % Maybe this is two chapters? The co-design could be one chapter, and the SPF itself could be another?

\minitoc

\section{Introduction}
To-do

\section{What Do We Mean When We Talk About `Diversity'?}
Talent Identification professionals frequently speak of diversity as a selection desideratum. This is variously used in the contexts of selected cohorts and selected individuals. Several definitions of diversity exist in the literature, but none encapsulate what TI professionals mean when they talk about diversity. We conduct a series of interviews with TI professionals investigating what they mean by "diversity", how they measure diversity, and why diversity is important. We conduct a thematic analysis of these interviews and report several themes related to the meaning of diversity. Finally, we discuss the significance of these insights and their applications to the development of purpose-built decision assistants designed to help TI professionals make diversity-conscious selection decisions.

(To-do)

\section{A Possibility Frontier Approach to Diverse Talent Selection} % If this gets too long, split this out into its own chapter.
Organisations (e.g., talent investment programs, schools, firms) are perennially interested in selecting cohorts of talented people. And organizations are increasingly interested in selecting diverse cohorts. Except in trivial cases, measuring the tradeoff between cohort diversity and talent is computationally difficult. Thus, organizations are presently unable to make Pareto-efficient decisions about these tradeoffs. We build on disparate understandings of diversity and introduce an algorithm that approximates the upper bound on cohort talent and diversity, as measured by one of a variety of target functions capturing different desiderata. We call this object the selection possibility frontier (SPF). We then use the SPF to assess the efficiency of the selection of a talent investment program run in 2021, 2022, and 2023. We show that, in the 2021 and 2022 cycles, the program selected cohorts of finalists that could have been better along both diversity and talent dimensions (i.e., considering only these dimensions as we subsequently calculated them, they are Pareto-inferior cohorts). But, when given access to our approximation of the SPF in the 2023 cycle, the program adjusted decisions and selected a cohort on the SPF.

(To-do)

\section{Co-designing Interpretable, Artificially Intelligent Decision Support Tools for Understanding Diversity}
The SPF represents one approach to resolving some of the concerns TI professionals have when discussing diversity. We interview TI professionals to understand what they would desire from tools designed to help them better address diversity concerns in selection. Furthermore, in addition to adapting the SPF methodology into a prototype designed to address specific TI professional concerns, we develop two other prototypes based on these interviews. Then, we run a scenario speed-dating exercise in order to understand the strengths and weaknesses of each approach. We evaluate each approach by how well it solves the problem it is intended to solve. Finally, we adapt these insights into a series of guidelines for how to design tools 

(To-do) % We can break down 4 and 5 by either TI problem domain or post-hoc vs. intrinsic. You can say there's good reasons to favor intrinsically interpretable models for both, but industry needs / existing work suggests post-hoc is necessary in that field
% \begin{savequote}[8cm]
% Alles Gescheite ist schon gedacht worden.\\
% Man muss nur versuchen, es noch einmal zu denken.

% All intelligent thoughts have already been thought;\\
% what is necessary is only to try to think them again.
%   \qauthor{--- Johann Wolfgang von Goethe \cite{von_goethe_wilhelm_1829}}
% \end{savequote}

\chapter{\label{ch:6-discussion}Discussion}

\minitoc

\section{Implications and Recommendations}
Here we analyse the implications of using IAIDSTs in talent identification. We first look at the strengths and limitations of post-hoc explanations and prescribe appropriate use cases for these types of IAIDSTs. Second, we examine the strengths and weaknesses of intrinsically interpretable models, including purpose-built models, and make similar prescriptions. Finally, we distil our insights in purpose-building IAIDSTs for TI professionals and extrapolate to general recommendations.

\section{Limitations and Future Work}
We note the limitations of this research method. We examine first limitations of individual case studies and note how other methodologies might address these limitations. We then raise concerns about our research's external validity, especially with regard to our recommendations and prescriptions. We address these concerns and note ways in which this research should not be used. Finally, we highlight areas in which researchers may continue this work or expand upon it.

\section{Conclusion}
To-do

% new chapter structure:
% intro
% % additional reasons to care about things xai overlooks because of the application domain
% methods
% limitations
% xai
% iai codesign
% implementing spf
% student ai use
% discussion


%%%%% APPENDICES
% Starts lettered appendices, adds a heading in table of contents, and adds a
%    page that just says "Appendices" to signal the end of your main text.
\startappendices
% Add or remove any appendices you'd like here:
\include{text/appendix-1}


%%%%% REFERENCES

% JEM: Quote for the top of references (just like a chapter quote if you're using them).  Comment to skip.
% \begin{savequote}[8cm]
% The first kind of intellectual and artistic personality belongs to the hedgehogs, the second to the foxes \dots
%   \qauthor{--- Sir Isaiah Berlin \cite{berlin_hedgehog_2013}}
% \end{savequote}

\setlength{\baselineskip}{0pt} % JEM: Single-space References

{\renewcommand*\MakeUppercase[1]{#1}%
\printbibliography[heading=bibintoc,title={\bibtitle}]}


\end{document}
